\documentclass[../main.tex]{subfiles}
\begin{document}
%==========================================TIÊU ĐỀ TẬP========================================
\begin{table}[h]
    \begin{adjustwidth}{-1cm}{}
        \begin{tabular}{>{\centering\arraybackslash}p{6cm}|>{\raggedright\arraybackslash}p{14cm}}
            \multirow{5}{6cm}
            % Cột bên trái: ÉPISODE và số tập
            {\\[-40pt]
                \flaregothic\fontsize{20pt}{20pt}\selectfont \centering
                \color{eptype}HISTOIRE\color{black}\\
                \gangofthree\fontsize{72pt}{72pt}\selectfont
                \color{epnum}10
                \\[18pt]
                % \flaregothic\fontsize{14pt}{14pt}\selectfont
                % \color{epattr}BIENVENUE, \uline{\color{Banpire} Mel} !
                \flaregothic\fontsize{18pt}{18pt}\selectfont
                \color{epattr} FIN DE TOME !
                \color{black}
            }
             &
            % Dòng thứ nhất: tên tập tiếng Nhật
            {\fotskip\fontsize{18pt}{18pt}\selectfont \ruby{新}{しん}\ruby{年}{ねん}、\ruby{古}{ふる}き\ruby{良}{よ}き\ruby{日}{ひ}々に\ruby{別}{わか}れを}
            \\[6pt]
            % Dòng thứ hai: tên tập tiếng Anh
             & {\cafeteria\fontsize{18pt}{18pt}\selectfont A Somber Way to End the Tet Holiday}  \\[4pt]
            % Dòng thứ ba: tên tập tiếng Việt
             & {\vnmsans\fontsize{18pt}{18pt}\selectfont Đám tang cuối lễ Tết}                   \\[8pt]
            % Dòng thứ tư: Nhân vật xuất hiện
             & {\brian{\fontsize{14pt}{14pt}\selectfont Track used: Lunar Zoo Year - Zombotany}}
        \end{tabular}
    \end{adjustwidth}
\end{table}
%===========================================NỘI DUNG==========================================
\color{black}

\pov{Hikawa Kuon}

\dots

\dots Tôi đã từng nói rằng, tôi muốn chúc tất cả mọi người có một năm mới bình an và hạnh phúc chứ?

Thực ra thì\dots

\dots ước gì tôi có thể nói như vậy với gia đình Hikawa chúng tôi.

\il

Chúng tôi đã ăn xong bữa cơm Rằm Tháng Giêng ở nhà ngoại tràn đầy ấm áp và nhẹ
nhàng. Bầu không khí vui vẻ của ngày Tết vẫn chưa tắt hẳn, chúng tôi vẫn còn chúc nhau
một năm mới tràn đầy sức khỏe và lợi lộc\dots

\dots cho đến khi tôi nhận được tin không hay, ngay tại gia đình nhà tôi.

Mười một giờ đêm. Tôi vẫn nhớ rõ lúc đó, mọi thứ bỗng trở nên căng thẳng. Dưới căn
phòng mà ông bà nội hay ngủ chung, ánh đèn vẫn còn sáng trưng. Mọi khi, đây đã là giờ
mà chúng tôi đã phải đi ngủ cho ngày hôm sau\dots

\dots nhưng rồi, một biến cố ập tới. Bác Hijaki-san và bố tôi đang cố gắng hết sức để giành
lại chút hy vọng nào cho ông nội, người đang nằm liệt giường.

Máy đo nhịp tim cứ bíp, bíp\dots\ nó vang lên từng tiếng như thể đang thử thách chúng tôi.
Các bác đứng xung quanh, tất cả nín thở, hy vọng một phép màu nào đó sẽ xảy ra.

Tuy nhiên, ông nội tôi dù gì cũng đã hơn tám mươi tuổi rồi, và ông cũng đã phải trải qua
vài lần đột quỵ, nên sức khỏe ông giờ đã giảm sút rõ rệt. Cả gia đình chúng tôi đã chuẩn
bị cho kịch bản không ai mong muốn này, nhưng vào cận kệ những giây phút cuối cùng
của ngày Tết lại là điều mà không ai trong chúng tôi lường trước được.

Không thể ngồi yên, tôi chạy xuống hỏi bố: ``{\color{red}Có chuyện gì thế bố?}'' với hai tay đang gắng
hết sức để ép ngực ông, cố gắng giành giật lấy sự sống rất mỏng manh kia.

Và rồi, câu trả lời của bố tôi khiến tôi không khỏi sững sờ: ``\uline{{\color{red}Ông mất rồi!}}''

\dots Dù tôi và ông không có mối quan hệ quá sâu nặng, sự mất mát của gia đình vào cái đêm
ấy vẫn khiến tôi trằn trọc, không thể ngủ được. Những gì tôi cảm nhận lúc đó\dots\ không
dễ dàng để diễn tả. Một phần của tôi muốn cố gắng tìm lý do, để không cảm thấy trống
vắng, nhưng tôi không thể.

Điều gì đến sau đó cũng đã phải đến. Khi tôi lên nhà lại được một lúc lâu, máy đo nhịp
tim ông nội giờ chỉ còn là một thẳng dài, kèm theo một tiếng *bíp* kéo dài dường như vô
tận.

Ông nội tôi đã qua đời vào đêm hôm đó, ngay trước khi ngày Tết thực sự kết thúc.

\ct

Đám tang sáng ngày hôm sau diễn ra trong một bầu không khí trầm lắng. Đại gia đình từ
các bác đến các con cháu đều có mặt, nhưng trong mắt ai cũng hiện lên nỗi buồn không
thể che giấu. Có vài người đeo khẩu trang, không chỉ vì đại dịch mà còn là một cách để
giấu đi cảm xúc.

Tuy nhiên, dù có làm thế nào, những khuôn mặt buồn rầu, mệt mỏi vẫn cứ lộ ra. Không
ai có thể phủ nhận rằng không khí trong nhà lúc đó rất nặng nề.

Tôi vốn dĩ là một người luôn giữ vững tinh thần. Nhưng kể cả vậy, tôi vẫn không cảm
thấy có chút gì đó thiêu thiếu bên trong lòng của mình, dù rằng ông nội và tôi không có
một mối quan hệ bền chặt.

Tôi đã xin nghỉ thêm ba ngày để có thể cùng dự đám tang của ông. Tanigo-san và A-san
cũng đã liên lạc với tôi, đồng thời gửi lời chia buồn và động viên từ xa. Dù không thể
ở cạnh, những lời động viên ấy vẫn mang lại cho tôi sự an ủi, dù biết rằng khoảng cách
không thể xóa nhòa nỗi đau lúc này.

Khi chúng tôi đi từ nhà đến khu mộ, những bước chân của gia đình nặng trĩu, mỗi người
đều chìm trong suy nghĩ của riêng mình. Dù không ai lên tiếng, nhưng sự im lặng giữa
chúng tôi như thể là một phần của nghi lễ này. Khi đến nơi, mắt tôi nhìn thấy một cảnh
tượng lạ thường, ngay tại khu đất mà chúng tôi từng đến mùng Ba Tết.

Có một chiếc huyệt đã được đào sẵn, nằm ở một góc đất trong khu mộ của gia đình. Tôi
không rõ từ lúc nào và ai đã đào nó. Chỉ biết rằng, khi nhìn thấy chiếc huyệt đó, cảm giác
bất an bỗng chốc tràn ngập trong lòng. Không ai trong gia đình tôi nhận ra nó, và điều đó
khiến bầu không khí thêm phần nặng nề.

Như thể rằng, có ai đó đã tiên tri tới một điều xảy đến với gia đình chúng tôi vậy.
Bố tôi nhíu mày, ánh mắt dò xét, rồi quay sang hỏi bác Hijaki-san: ``{\color{red}Cái này là sao vậy hả anh? Ai đã đào nó lên?}''

Bác lắc đầu, trả lời: ``{\color{red}Tao cũng không biết. Chắc chắn là không phải của gia đình mình\dots\ Từ lúc tao về đây mùng Năm để thăm, tao còn chưa thấy cái huyệt này cơ.}''

Mọi người im lặng nhìn nhau, không ai dám lên tiếng, chỉ có tiếng bước chân trên đất và
những tiếng thở dài.

Mẹ tôi cúi xuống, đưa tay sờ vào lớp đất còn mới, cảm giác lạnh lẽo ập đến. ``Nó\dots\ không
giống với những huyệt chúng ta đã đào trước đây.'' Giọng bà run lên một chút, như thể
sự không bình thường của chiếc huyệt này đã tác động đến tâm lý của bà. Một cảm giác
gì đó vừa buồn, vừa sợ đến gáy lưng.

Dù không có lời giải thích, nhưng sự xuất hiện bất ngờ của chiếc huyệt cứ đeo bám chúng
tôi cho đến tận lúc đưa hũ tro của ông vào nơi an nghỉ. Không ai dám nói ra, nhưng tôi
cảm thấy như có điều gì đó bất thường đang diễn ra trong gia đình mình, không chỉ với
cái huyệt lạ mà còn với những cảm xúc lạ lẫm mà nó mang lại.

Cái huyệt đó, có lẽ sẽ là câu hỏi không lời đáp trong gia đình tôi mãi mãi. Và dù ai đó có
muốn lý giải, nó vẫn sẽ là một phần của ký ức, như một dấu hiệu mà thời gian sẽ chẳng
thể xóa nhòa được.

Lễ truy điệu diễn ra trong yên lặng. Mọi người không nói nhiều, chỉ lặng lẽ đi từ nhà của
cậu đến đài hóa thân hỏa táng, nơi hũ tro của ông được đặt xuống. Sau đó, hũ tro được đưa
về khu mộ, nơi mà chính tôi, Miko-chan, Fubu-chan và Subaru-chan cùng Shion-chan đã đến vào mùng Ba Tết. Chúng tôi đã đi từ nhà đến mảnh đất ấy, nhưng hôm nay, nơi đây
lại trở thành dấu chấm hết cho một vòng đời.

Khi chiếc hũ tro của ông được đặt xuống cái huyệt bí ẩn kia và những lớp đất được lấp
lại, tôi cảm thấy như một phần của gia đình mình đã thực sự rời xa. Khoảnh khắc đó thật
khó để diễn tả bằng lời, nhưng có thể nói, đó là một cái kết buồn cho gia đình Hikawa nói
riêng và gia đình tôi nói chung. Một kết thúc không thể nào thay đổi.

Từ đây trở đi, gia đình tôi sẽ không còn đủ đầy như trước nữa. Mọi người sẽ bận rộn hơn,
công việc sẽ khiến chúng tôi ít gặp nhau hơn, và đương nhiên, sẽ không còn những dịp
tụ họp đông đủ như ngày Tết nữa. Dù mỗi lần gặp mặt vẫn sẽ có ý nghĩa, nhưng sự thiếu
vắng của ông khiến không khí luôn có sự trống vắng lạ lùng.

Nhưng trước mắt, năm đó, cả nhà Hikawa chúng tôi xem như ``mất Tết''. Mọi niềm vui,
mọi lễ hội mà tôi và các \hll\ đã tạo nên trước đó, đều không thể làm mờ đi sự mất
mát mà gia đình tôi vừa trải qua. Tết Nguyên Đán năm đó thật sự không còn là dịp để vui
mừng như những năm trước nữa, mà chỉ còn lại sự tĩnh lặng, nặng trĩu những suy tư và
tiếc nuối.

\pov{Murasaki Shion}

Tôi ngồi một mình trong căn phòng tối, ánh sáng duy nhất là từ quả cầu ma thuật phát ra
một luồng sáng mờ nhạt.

Tại sao tôi lại nhìn chằm chằm vào quả bóng này sao? Bằng một cách nào đó, cái cảm
giác bất an đó trào dâng lên đầy kỳ lạ, kể từ đêm hôm qua. Tôi đã không ngủ được.

Tôi đã thức suốt đêm để theo dõi gia đình nhà Kuon-senpai, và những gì tôi thấy đã làm
tôi giật mình. Một thứ gì đó mà thật sự không ai mong muốn.

Ông nội của anh ấy đã ra đi, đúng vào những giờ cuối cùng của ngày Tết.

Buổi sáng theo sau đó, tôi nhìn thấy hình ảnh gia đình của anh ấy đang đứng lặng lẽ bên
mộ ông nội, một cảnh tượng trang trọng nhưng cũng đầy ảm đạm. Những khuôn mặt quen
thuộc, những bóng hình mà tôi đã biết từ lâu, giờ đây mang một vẻ đau buồn khó diễn tả.

Tôi có thể cảm nhận được sự tĩnh lặng trong không khí, như thể một phần nào đó của họ
đã ra đi mãi mãi. Lễ truy điệu đang diễn ra, nhưng cái nhìn của tôi không chỉ dừng lại ở
đó. Một sự hiện diện mà tôi không thể chối bỏ, sự hiện diện của chiếc huyệt mà tôi đã
đào.

Ngay khi chiếc huyệt ấy xuất hiện trong tầm mắt tôi qua quả cầu, một cảm giác lạnh lẽo
dâng trào trong lòng. Cả người tôi như đông cứng lại, và tôi không thể rời mắt khỏi nó.
Chẳng biết vì sao, nhưng tôi cảm thấy như mình đã biết điều này sẽ xảy ra từ lâu rồi.

Tôi đã cố gắng thuyết phục bản thân rằng đó chỉ là một sự tình cờ, một sự kiện vô tình,
nhưng giờ đây, khi thấy nó - cái huyệt mà gia đình anh ấy đã dùng để chôn cất hũ tro ông
nội đi - thực sự xuất hiện trước mặt tôi, nỗi sợ hãi sâu thẳm trong lòng tôi đã thành sự thật.

Không lẽ rằng tôi đã tiên tri cho một điều gì đó không lành với gia đình của anh ấy?

``\textit{L-Là mình đã làm vậy sao?}'' tôi thì thầm, giọng nói của mình nhỏ đến mức như thể
không thể thoát ra khỏi chính tôi. Tôi đã không thể ngừng nghĩ về chiếc huyệt đó, không
thể ngừng tự hỏi liệu nó có phải là dấu hiệu của một điều gì đó sắp xảy đến. Cảm giác
tội lỗi không thể dứt bỏ cứ bám riết lấy tôi, như thể chính tôi đã vô tình mở ra một con
đường mà không ai có thể quay lại.

Và rồi, đám tang của nhà anh ấy đã minh chứng cho cái điềm báo không lành ấy.
Tôi không thể dối lòng mình nữa. Đó không phải là một sự tình cờ, không phải là một sự
việc ngẫu nhiên. Tôi đã tạo ra nó. Tôi đã đào chiếc huyệt đó, và giờ nó đang chứng minh
rằng những điều tồi tệ mà tôi lo sợ cuối cùng cũng đãn.

Tôi rũ xuống, ngồi thụp trên sàn nhà. Trái tim tôi như bị bóp nghẹt, từng nhịp đập chậm
lại trong sự căng thẳng. Tôi cảm thấy mình thật sự cô độc, thật sự lạc lõng trong những
cảm xúc hỗn độn. Không ai ở đây để chia sẻ nỗi sợ hãi này với tôi. Và tôi không thể
không tự hỏi, liệu có phải tôi đã gây ra nỗi đau cho gia đình Kuon-senpai, cho ông nội
anh ấy, và cả những người tôi luôn coi là bạn bè?

Sự im lặng của căn phòng càng làm cho tâm trí tôi thêm tăm tối và rối bời. Tôi không
dám nhìn thẳng vào hình ảnh trong quả cầu nữa, tôi không thể chịu nổi cái cảm giác ấy
thêm một giây nào nữa. Dù có muốn, tôi cũng không thể ngừng nghĩ về chiếc huyệt - và
những gì nó mang đến.

Ngày Tết đã kết thúc, nhưng đối với tôi, cảm giác như thời gian đã đứng im. Dù những
ngày vui vẻ, những khoảnh khắc ngập tràn tiếng cười vẫn đang lặng lẽ trôi qua trong ký
ức, nhưng nỗi ám ảnh từ chiếc huyệt hôm đó cứ vương vấn, không sao thoát ra được. Tôi
không biết phải làm gì với cảm xúc của mình, với nỗi tội lỗi mà tôi đã tạo ra.

Tôi ngồi trong phòng, ánh sáng mờ ảo từ quả cầu ma thuật vẫn còn chiếu sáng mặt tôi.
Lúc này, tôi chỉ muốn trốn tránh mọi thứ. Trốn tránh cảm giác tội lỗi, trốn tránh việc phải
đối diện với những người mà tôi đã lỡ làm tổn thương. Nhưng lại có một phần trong tôi
biết rằng, dù có trốn tránh, cuối cùng cũng phải đối mặt với tất cả.

Cảnh tượng gia đình của anh, những gương mặt buồn bã trong lễ truy điệu, vẫn hiện rõ
trong đầu tôi. Tôi biết, dù tôi không ở đó, họ cũng sẽ không bao giờ quên những gì đã
xảy ra. Và tôi, trong một cách nào đó, sẽ phải sống với nỗi sợ đó mãi mãi.

Dường như với tôi, Tết không còn là dịp vui vẻ nữa. Đối với tôi, ngày Tết này đã đánh
dấu một sự kết thúc\dots\ một kết thúc không mong muốn. Dẫu biết rằng những cái Tết sau
có thể sẽ tốt đẹp hơn so với năm nay, nhưng liệu rằng nó sẽ còn tươi đẹp và vui vẻ như
những ngày trước thay không?

Tôi không biết liệu có thể sửa chữa những sai lầm này không, liệu tôi có thể làm gì để bù
đắp cho gia đình Kuon-senpai, nhưng tất cả những gì tôi có thể làm lúc này là ngồi im,
chấp nhận những gì đã xảy ra.

\pov{Hikawa Kuon}

Ngày Tết năm đó đã không như tôi mong đợi. Dù đã có sự chuẩn bị, dù đã có những
khoảnh khắc vui vẻ khi cả gia đình cùng nhau đón Tết, nhưng nó không thể nào xóa nhòa đi những gì đã xảy ra.

Lễ truy điệu sáng hôm sau là một khoảnh khắc không thể nào quên. Mọi người trong gia
đình đứng bên cạnh mộ ông nội, sự tĩnh lặng trong không khí làm tôi cảm thấy như cả
thế giới này đang thu nhỏ lại. Tôi không thể không nghĩ về chiếc huyệt mà ai đó đã đào.
Điều đó cứ ám ảnh tôi, dù tôi có muốn nó không xuất hiện trong đầu mình.

Chắc chắn không ai trong chúng tôi có thể hiểu được tại sao nó lại ở đó. Không ai hiểu
tại sao nó lại xuất hiện vào lúc này. Và dù có muốn đổ lỗi cho bất cứ ai, tôi vẫn không thể
phủ nhận rằng mọi thứ dường như đã sẵn sàng. Tôi không thể xua đi được cảm giác rằng
tất cả chỉ là một chuỗi sự kiện dẫn đến khoảnh khắc này.

Nhìn ông nội lần cuối trước khi hỏa táng, tôi cảm thấy một khoảng trống vô tận. Tết này,
gia đình chúng tôi không còn đủ đầy nữa. Mọi thứ đã thay đổi, và có lẽ sẽ không bao giờ
như xưa. Để lại trong lòng tôi một nỗi buồn sâu sắc, một sự thiếu vắng mà không thể bù
đắp.

Tôi không thể quay lại những khoảnh khắc vui vẻ của những ngày trước Tết. Mọi thứ đã
thay đổi quá nhanh, và tôi biết, dù có muốn thế nào đi nữa, Tết năm nay sẽ mãi là một ký
ức đắng cay trong lòng mình, và trong lòng của tất cả những người đã tham dự vào đám
tang ngày hôm đó.

\end{document}